% !TEX root = ../main.tex

\section{Analytical Solution for a Dipole} \label{sec:dipole_derivation}

Dipoles or loading sources are sounds due to forces acting on the fluid, equivalent to body forces and momentum supply. It is also equivalent to two equal monopoles with opposite signs at an infinitesimally small distance. The dipole is described as the following,

\begin{equation}
    p'(\bm{x}, t)=-\frac{\partial}{\partial x_i} \left[\frac{f_i(\tau)}{4\pi r(1-M_r)}\right]_{\tau *}
    \label{eq:dipole_convolved}
\end{equation}

In the dipole description above, the source forcing term $f_i$ is not an arbitrary vector. To analytically model a dipole, the mathematical equivalence between monopoles and dipoles must be established. Sound pressure generated by a monopole source is defined as,

\begin{equation*}
    p'_{mono}(\bm{x}, t) = \frac{\partial}{\partial t} \left( \frac{Q(\tau)}{4\pi r|1-M_r|} \right)
\end{equation*}

By describing a dipole as two monopoles of equal strength and opposite signs separated by an infinitely small distance $\delta$, the dipole pressure is a superposition of the two monopoles,

\begin{equation*}
    p'_{dipole}=p'_{mono}(\bm{x},t)-p'_{mono}(\bm{x}-\delta, t)
\end{equation*}

As $\delta$ tends to zero, this distance can be expressed as a first-order Taylor series expansion,

\begin{equation*}
    p'_{dipole}(\bm{x},t) \approx -\delta \frac{\partial}{\partial x_i} \frac{\partial}{\partial t} \left[ \frac{Q(\tau)}{4\pi r|1-M_r|} 
    \right] 
\end{equation*}

Hence, the dipole source strength $f_i(t)$ can be defined as,

\begin{equation}
    f_i(t)=4\pi r|1-M_r| \delta \frac{\partial}{\partial t} \left(\frac{Q(t)}{4\pi r |1-M_r|} \right)
    \label{eq:dipole_source}
\end{equation}

by defining the forcing term in this term and evaluating it using central differencing, the analytical code captures the expansion of this forcing term without requiring an analytical expansion.

\subsection{Dipole Pressure Equation} \label{subsec:dipole_p}

Using the definition described in Equation \ref{eq:dipole_convolved}, replacing the $\partial / \partial x_i$ term as outlined in Appendix~\ref{app:ddx_replace}, and using the steps from monopole pressure, the pressure for an analytical dipole is obtained,

\begin{equation}
\begin{split}
    p'(\bm{x}, t) &= \frac{(\partial f_i(\tau)/\partial \tau) r_i}{4\pi cr^2(1-M_r)^2}
    + \frac{f_i(\tau)r_i}{4\pi r^3 |1-M_r|^3}
    - \frac{f_i(\tau)M_i}{4\pi r^2 |1-M_r|^2} \\
    &\phantom{=} + \left( \frac{\partial M}{\partial \tau}\right)_r \frac{f_i(\tau)r_i}{4\pi cr^2|1-M_r|^3}
    + \frac{f_i(\tau)r_iM_j^2}{4\pi r^3 |1-M_r|^3}
    \label{eq:dipole_pressure}
\end{split}
\end{equation}

where the first and fourth terms on the right-hand side are the far-field propagation due to changes in the source forcing term and Mach number, respectively; while the second, third, and fifth terms are the near-field propagation terms. The far-field terms are the primary reason for the noise propagation in rotors such as helicopter blades, propellers, and open rotor engines; as the blades rotate and push the air surrounding it, the local forcing strength and Mach number changes with time.

\subsection{Dipole Velocity Equation} \label{subsec:dipole_u}
To analytically derive the velocity for a dipole, explicitly integrating the momentum equation like with the monopole is unviable, as this approach will return too many terms. Instead, limited solutions are obtained by applying some physical assumptions. The velocity derivations assume either a zero relative Mach number or a plane wave condition. An alternative method that uses the superposition of two equal monopoles of opposite signs will also be presented.

\subsubsection{Dipole Velocity Equation: Stationary Assumption} \label{subsubsec:dipole_u1}

Assuming inviscid and no mean flow in the momentum equation, the pressure gradient is,

\begin{equation*}
    -\frac{\partial p}{\partial x_i} = 
    \frac{\partial^2}{\partial x_i \partial x_j} \left[ 
    \frac{f_i(\tau)}{4\pi r (1-M_r)}
    \right]_{\tau*}
\end{equation*}

Hence, the velocity is defined as,

\begin{equation*}
    u_i' = \frac{1}{4\pi\rho_0} \frac{\partial^2}{\partial x_i \partial x_j} \int_{-\infty}^t \frac{f_i(\tau)}{r(1-M_r)}dt
\end{equation*}

substituting the integrand to use the time at the observer, $\tau$, and assuming zero relative velocity; therefore, $dr/dt=0$ and $M_r=0$. These assumptions are what causes this assumption to only be valid when the source is stationary; such that a co-translation with the observer will still violate this assumption. 

To resolve the integral above, let $F_i(\tau) = \int f_i(\tau) d\tau \implies f_i(\tau) = \dot{F}_i(\tau)$, the analytical solution for dipole velocity that is valid for cases with zero relative Mach number is obtained,

\begin{equation}
    u_i'=\frac{1}{4\pi\rho_0} \left( 
    \frac{F_j(3\textbf{n}_i \textbf{n}_j-\delta_{ij})}{r^3} +
    \frac{f_j(3\textbf{n}_i \textbf{n}_j - \delta_{ij})}{cr^2} +
    \frac{\dot{f_j}(\textbf{n}_i \textbf{n}_j)}{c^2r}
    \right)
    \label{eq:dipole_u1}
\end{equation}

where the first, second, and third terms on the right hand side are the near-, intermediate-, and far-field velocity propagation, respectively. 

\subsubsection{Dipole Velocity Equation: Plane Wave Assumption} \label{subsubsec:dipole_u2}

The plane wave assumption allows relative Mach number between the source and observer, provided the control surface is in far-field that it encapsulates the source's entire wavelength. In the far-field, the pressure and velocity field perturbation components are in-phase, which is the reason of this requirement. The far-field is defined as,

\begin{equation}
    kR \gg 1
    \label{eq:far_field_def}
\end{equation}

% where $\hat{r}$ is the control surface radius. 
where $k$ is the wave number. Physically, this means the acoustic wavelength must be smaller than the control surface diameter, allowing the entire wavelength to be fully captured by the permeable sphere. Using the derivations in Appendix~\ref{app:planeWave}, the velocity equation is obtained,

\begin{equation}
    u_i' = \frac{\textbf{n}_i}{\rho_0c}p'(\bm{x},t)
    \label{eq:dipole_u2}
\end{equation}

hence showing that the velocity perturbation component is in-phase with the pressure.

\subsubsection{Dual Monopole} \label{subsubsec:dual_mono}

An alternative approach to validate dipole sources is the dual monopole method. This approach entirely avoids the assumptions that were made in the previous derivations and instead uses two monopoles to create a dipole. This bypasses the limitations imposed by the assumptions; therefore, this method will remain accurate when used on cases that would violate the above assumptions.. However, to ensure that the dual monopole approach itself is valid, it needs to be compared against the analytical solutions above.

Using the definition of a dipole, this method uses two equal monopoles of opposing signs, separated by a small distance $\delta$, which in this case is defined as,

\begin{equation}
    k\delta \ll 1
    \label{eq:delta_distance_def}
\end{equation}

where $k$ is the wave number. To simplify the validation process, it is that the distance between the two monopoles is $\delta=0.01/k$, which still validates this method. Another advantage of this approach is that it allows the prediction of the directivity patterns of the dipole, the results of which will also be presented in the next section.